\section{The constraint space and the theorem}
\label{sec:phase}

We use Lorentzian signature \((-+++)\). Fix \(M>0\) and
\(0<\mathfrak a<M\), where \(\mathfrak a\) denotes the Kerr rotation
parameter. The letter \(a\) will later denote a double-null metric
coefficient, unrelated to this parameter.

Let \((\Sigma,h_*,K_*)\), \(\Sigma\simeq\R\times\mathbb S^2\), be the
complete signed Killing-time-zero Kerr bridge through the bifurcation
sphere, with the future-normal convention of
\cite[\CompanionRef{foundation}{sec:introduction} and \CompanionRef{foundation}{sec:comparison}]{Foundation}. Fix its identification with \(\Sigma\),
asymptotically Cartesian charts on both ends, and a smooth function
\(\rho\ge1\) equal to the end radius outside a compact set.
Write \(\nabla_*\) and \(|\cdot|_*\) for the reference connection and
tensor norm.

\begin{definition}[Weighted constraint space]
\label{phase:space}
For smooth symmetric two-tensors \(f,j\) set
\begin{equation}
\label{phase:seminorms}
 p_m(f,j)=\sum_{i=0}^m
 \left\{\sup_\Sigma \rho^{1+i}|\nabla_*^if|_*
       +\sup_\Sigma \rho^{2+i}|\nabla_*^ij|_*\right\},
 \qquad m=0,1,\ldots.
\end{equation}
Let \(\mathcal D\) consist of smooth pairs \(d=(h,K)\) such that \(h\)
is positive definite and complete, every \(p_m(h-h_*,K-K_*)\) is finite,
and
\begin{equation}
\label{phase:constraints}
 R(h)+(\tr_hK)^2-|K|_h^2=0,\qquad
 \operatorname{div}_hK-\dd(\tr_hK)=0.
\end{equation}
The topology on \(\mathcal D\) is the relative topology of the seminorms
\eqref{phase:seminorms}. For \(\eps>0\) let
\[
 \mathcal U_\eps=\{d\in\mathcal D:p_{10}(d-d_*)<\eps\},
 \qquad d_*=(h_*,K_*).
\]
\end{definition}

These are differentiated \(O(r^{-1})\) and \(O(r^{-2})\) end conditions.
There is no parity, fixed leading coefficient, fixed charge, or symmetry
restriction. In particular, we do not require every possible
asymptotic charge integral to exist for every allowed tail.
We use representatives in the fixed spatial gauge, rather than a
quotient by diffeomorphisms.

The choice of equivalent end charts, weights and reference geometry
does not change the topology. More precisely, suppose the radii are
comparable, the transition Jacobian and its inverse are symbols of
order zero, and the transformed reference geometry has the same
asymptotic class. The tensor chain rule and the formula for changing
connections bound each seminorm in one choice by a finite sum of
seminorms in the other, in both directions. To retain the Cartesian
convention, the leading Euclidean metric is preserved up to the
allowed metric-weight error. A smooth reidentification carrying both
the data and the reference center therefore gives a homeomorphism.
This assertion does not identify seminorms under arbitrary
data-dependent gauge transformations.

For \(d\in\mathcal D\), let \(\mathcal M_d\) denote its full smooth
maximal globally hyperbolic vacuum development, with the prescribed
time orientation. We use smooth Cauchy existence, propagation of the
wave-gauge constraints, geometric uniqueness, and compact-subset
Cauchy stability in their usual smooth-data form; the required local
estimates are part of the Cauchy framework in
\cite[\CompanionRef{foundation}{sec:category}]{Foundation}.
The defining embedding property of the maximal development is
\cite[Definitions 2.1--2.2 and Theorems 2.6--2.8]{Sbierski2016}.
An application of that property below will always be preceded by a
proof that the attached spacetime has the \emph{entire} initial
manifold as a Cauchy surface.

\begin{definition}[Future \(C^1\) extension]
\label{phase:extension}
A future \(C^1\) extension of \(\mathcal M_d\) is a smooth,
time-orientation-preserving isometric embedding
\(\iota:\mathcal M_d\longrightarrow\widetilde{\mathcal M}\)
onto a proper open subset of a connected, time-oriented smooth
four-manifold with nondegenerate \(C^1\) Lorentzian metric
\(\widetilde g\), together with a future-directed timelike \(C^1\)
curve \(c:[0,1]\to\widetilde{\mathcal M}\) such that
\[
 c([0,1))\subset\iota(\mathcal M_d),\qquad
 c(1)\in\partial\iota(\mathcal M_d).
\]
Neither the vacuum equations nor global hyperbolicity are imposed
on the ambient spacetime.
\end{definition}

\begin{theorem}
\label{phase:main}
For each fixed \(M>0\) and \(0<\mathfrak a<M\), there is
\(\eps_0=\eps_0(M,\mathfrak a)>0\) such that
\(\mathcal U_{\eps_0}\) contains a dense \(G_\delta\) subset
\(\mathcal G\) for which every \(\mathcal M_d\), \(d\in\mathcal G\),
has no future \(C^1\) extension in the sense of
Definition~\ref{phase:extension}.
\end{theorem}

The open neighborhood and the residual subset in this statement
are different sets. The theorem does not assert that inextendibility
is open, or that every datum in \(\mathcal U_{\eps_0}\) has the
conclusion. Its constants need not be uniform as
\(\mathfrak a\to0\) or \(\mathfrak a\to M\).
All later choices of the base neighborhood use only the smallness of
\(p_{10}\); higher seminorms are finite but need not be uniformly small.
Appendix~\ref{phase:symmetry} constructs nearby data with trivial spatial
data isometry group and explains how nearby Kerr centers fit into the
same phase space. These observations are separate from the genericity
argument.

\begin{lemma}[Baire property]
\label{phase:baire}
The space \(\mathcal D\) is completely metrizable, and every
\(\mathcal U_\eps\) is a nonempty Baire space. At any point, finite
seminorm balls form a neighborhood base: every neighborhood contains
a set of the form
\[
 \{d'\in\mathcal D:p_m(d'-d)<r\}
\]
for some finite \(m\ge10\) and \(r>0\).
\end{lemma}

\begin{proof}
The vector space \(\mathcal E\) of tensor pairs with every seminorm
\eqref{phase:seminorms} finite is complete for the usual metric built
from the countable increasing seminorms. Indeed, a Cauchy sequence
converges smoothly on compact sets, and its weighted uniform
derivative limits agree with the derivatives of that local smooth
limit. Passing the uniform Cauchy bounds to the limit gives
convergence in each \(p_m\).

For any \(h=h_*+f\) with \(h>0\) and \(p_0(f,0)<\infty\), the
asymptotic decay makes \(h\) uniformly comparable with \(h_*\) on
the ends. Positivity on the remaining compact set gives the same
comparison there. Thus \(h\) is complete. The same positive lower
comparison bound shows that positivity is open in \(\mathcal E\).
Consequently the completeness requirement in Definition~\ref{phase:space}
is automatic on this positive open set.

The constraint map is continuous there into the space of smooth
scalar and one-form fields with the local smooth topology.
Its zero locus is relatively closed. An open subset of a completely
metrizable space is completely metrizable, as is a closed subset
of it. This proves the assertion for \(\mathcal D\), and then for
its open subsets \(\mathcal U_\eps\). They contain \(d_*\).
Finally the seminorms increase with \(m\), so finitely many
seminorm inequalities are implied by one sufficiently high-order
inequality with sufficiently small radius.
\end{proof}

\subsection{Mass normalization}
\label{phase:scaling}

The constant rescaling \(g\mapsto M^{-2}g\), written in dimensionless
bridge coordinates, sends the initial data to
\((M^{-2}h,M^{-1}K)\). It gives a homeomorphism between the
corresponding weighted constraint spaces, sends the reference
parameters to \((1,\mathfrak a/M)\), and preserves both maximality
and Definition~\ref{phase:extension}. We henceforth take \(M=1\)
and fix \(0<\mathfrak a<1\). Proving the theorem in this normalization
proves it for every fixed original center.

Choose \(\eps_0\) small enough for the background statements in
Section~\ref{sec:background}, and write
\(\mathcal U=\mathcal U_{\eps_0}\).
